% ctx-intro.tex
%
% Copyright (C) 2022--2026 José A. Navarro Ramón <janr.devel@gmail.com>
% Licencia del código GPLv2
% Licencia Creative Commons Recognition Non-Commercial Share-alike.
% (CC-BY-NC-SA)


\startcomponent ctx-intro
  \product tutorialLuaMetaFun
  
  \startchapter[title=Introducción]
    En este capítulo presentaremos unos pocos conceptos para abrir boca con
    \MetaFun.

    Empezaremos distinguiendo dos tipos de objetos gráficos de \MetaPost\ y de
    \MetaFun.
    \startitemize[n,packed,joinedup]
    \item {\bf \type{path}, o camino, trazo o línea}:
      Es como la silueta invisible que trazas en el aire con un dedo. No tiene
      color, grosor ni visibilidad. Es invisible hasta que lo dibujas con un comando.
    \item {\bf \type{picture}, imagen u objeto gráfico}:
      Es como el lienzo donde has aplicado pintura sobre la silueta \type{path}.
      Es un contenedor visual. Almacena el resultado de haber dibujado
      \type{draw} o rellenado \type{fill} un \type{path}, incluyendo sus colores,
      grosores de línea y textos. Hay objetos gráficos nativos de \MetaPost\ o
      \MetaFun, que se generan en el código nativo de estos sistemas. Pero hay otros
      que se originan a partir de un archivo gráfico con formato .svg, .pdf, .png,
      .jpg u otros. En este último caso se importan a \MetaPost / \MetaFun\ mediante
      la orden \type{externalfigure}.
    \stopitemize
    \blank[4em]
    Analicemos el siguiente dibujo y su código

    \blank[-2cm]
    % Definición de búfers
    \startbuffer[circulo_inicial]
      \startMPcode
        fill fullcircle;
      \stopMPcode
    \stopbuffer
    % Para código fuente
    \startbuffer[circulo_inicial_txt]
      \starttyping[option=METAFUN]
        \startMPcode
          fill fullcircle ;
        \stopMPcode
      \stoptyping
    \stopbuffer
    % .........................................................................
    % Impresión de búfers
    \startlinealignment[left]
      \framed[frame=off,location=middle]{%
        \getbuffer[circulo_inicial]}% <--- ¡El % aquí es vital!
      \hskip 3.0em % Aumentamos el espacio entre la imagen y el código
      \framed[frame=off,location=middle,align=right,width=fit]{%
        \switchtobodyfont[10.0pt]\typebuffer[circulo_inicial_txt]}
    \stoplinealignment

    Pero, ¿qué está pasando? ¿dónde está el dibujo? Si observas con atención verás
    un pequeño punto negro a la izquierda del código. Para explicar lo que está
    ocurriendo vamos a analizarlo.

  \startsection[title={Comandos como \color[blue]{\type{fill}}} y otros]

    El comando \type{fill} rellena de color ---negro por defecto--- la imagen o
    el dibujo (trazo) que le sigue en el código. En \MetaFun\ hay más comandos de
    dibujo y manipulación visual, como \type{draw}, \type{filldraw}, \type{undraw},
    \type{unfill}, \type{drawpoints}, \type{drawshadow}, \type{drawlabel}
    (o \type{label}), que iremos utilizando conforme avancemos. A continuación
    presentamos brevemente su utilidad:
    \startitemize[n, packed, joinedup]
      Comandos de combinación y relleno avanzado (\MetaPost):
      \startitemize[packed, joinedup]
      \item {\bf\type{fill}}:
        Rellena el trazo cerrado que se le pasa.
      \item {\bf\type{draw}}:
        Dibuja una línea (trazo), o una imagen creada en MetaFun o que proviene de
        un fichero gráfico (.svg, .png, jpg, etc.) que se le pasa al comando.
      \item {\bf\type{filldraw}}: Combina a la vez los dos comandos anteriores.
        Rellena el interior de un trazo cerrado y, al mismo tiempo, dibuja su
        contorno. Es más eficiente que realizar las dos operaciones por separado.
      \item {\bf\type{drawfill}}: Es el inverso conceptual de \type{drawfill}.
        En la práctica realiza la misma acción combinada, pero asegura que las
        propiedades del trazo exterior e interior se procesen juntas.
      \item {\bf \type{undraw}/\type{unfill}/\type{unfilldraw}}:
        Estos comando actúan como un \quotation{borrador eléctrico}. No pintan de
        color blanco; lo que hacen es \quotation{restar} o recortar un trazo o un
        área del objeto gráfico (\type{picture}) actual, dejando el fondo
        transparente.
      \stopitemize

    \item Comandos de \MetaFun\ (extensiones de \ConTeXt):
      \MetaFun\ introduce macros de alto nivel que simplifican el diseño editorial
      y de fondos (\emph{overlays}):
      \startitemize[packed, joinedup]
      \item {\bf\type{drafill} con tranparencias}:
        \MetaFun\ añade soporte nativo para transparencias mediante
        \type{withtransparency}.
      \item {\bf\type{drawpoints} / \type{drawcontrolpoints}}:
        Son comandos de depuración muy útiles. Dibujan de forma automática los
        puntos clave (nodos) o las manecillas de control \emph{Bézier} de un trazo
        para que veas cómo está construida la curva matemáticamente).
      \item {\bf \type{drawboundingbox}}: Dibuja un rectángulo exacto que delimita
        los bordes del marco invisible de un trazo o una imagen completa.
      \stopitemize

    \item Comandos para texto y sombras:
      \startitemize[packed, joinedup]
      \item {\bf\type{drawlabel} / \type{label}}:
        Coloca texto en coordenadas específicas del gráfico. \MetaFun\ añade
        \type{thelabel} para generar el texto como un objeto independiente antes
        de posicionarlo.
      \item {\bf \type{drawshadow}}:
        Es una macro exclusiva de \MetaFun\ que genera automáticamente una sombra
        difuminada o desplazada debajo de un trazo o figura.
      \stopitemize
    \stopitemize
  \stopsection

  \startsection[title=Función generadora \color[red]{\type{fullcircle}}]
    
    Ahora fijamos nuestra atención en el término \type{fullcircle} que sigue al
    comando. Es una función que \emph{genera} un trazo, camino o \emph{path} en
    inglés; en este caso es un camino cerrado que se puede rellenar con el comando
    \emph{\type{fill}}. Nótese que, por sí solo, \type{fullcircle} no dibuja nada,
    incluso da error, necesita un comando detrás. Es como un dedo en el aire
    dibujando un círculo.

    Un \emph{path}, como el que crea \type{fullcircle}, es solo geometría pura.
    Define una forma matemática (puntos, líneas y curvas), pero no tiene color,
    grosor ni visibilidad. Solo es visible cuando se dibuja o rellena.

    Entonces, ¿por qué es tan pequeño el círculo que se dibuja cuando lo rellenamos
    con el comando \type{fill}? Lo que ocurre es que, al igual que el color del
    círculo es negro por defecto, su diámetro es la unidad de longitud establecida
    en \MetaPost, que por defecto es \type{1 bp} o \quotation{punto grande},
    \emph{big point} en inglés.

    Es importante entender que esta función define la geometría, el tamaño y la
    posición del centro del círculo; la forma es un círculo, el tamaño es una unidad
    \MetaPost\ y su centro está en el origen (0,0). El comando \type{fill} solo
    ejecuta la acción de rellenar lo que genera \type{fullcircle}.

    \startsubsection[title=Unidad \color[black]{\type{bp}}]
    
      Esta unidad se define como 1/72 de pulgada, (0.35278 cm aproximadamente).
      Como se puede observar, su tamaño es muy pequeño.
    
    \stopsubsection

    \startsubsection[title=Operador secundario \color[darkgreen]{\type{scaled}}]
      Para poder modificar el tamaño del círculo se suele utilizar la función
      modificadora \type{scaled}.
      En la jerga interna de MetaPost/MetaFun, se clasifica como un
      \type{secondarydef} (operador secundario) o \emph{binary operator}.
      Funciona exactamente como un operador matemático (como el signo de multiplicar
      \type{*}), pero diseñado específicamente para transformar objetos geométricos o
      numéricos. Este operador modifica las dos dimensiones, horizontal y vertical
      en la misma proporción.

      El factor de escalado se puede escribir directamente o entre paréntesis.
      En el siguiente ejemplo rellenaremos un cículo de \unit{1 cm} de diámetro
      con el color negro por defecto y con su centro en el origen:
      
      \startlinealignment[left]
        \framed[frame=off,location=middle]{%
          \startMPcode
            fill fullcircle scaled 1cm ;
          \stopMPcode}
        \hskip -2.0em % Disminuimos el espacio entre la imagen y el código
        \framed[frame=off,location=middle,align=right,width=fit]{%
          \switchtobodyfont[11.0pt]
          \starttyping[option=METAFUN]
            \startMPcode
              fill fullcircle scaled 1cm ;
            \stopMPcode
          \stoptyping}
      \stoplinealignment

      ¿Y si no especificamos la unidad? Entonces \MetaFun\ utiliza la unidad por
      defecto de \MetaPost, que es \type{bp}. En el ejemplo siguiente, estamos
      rellenando un círculo de 20bp de diámetro (20 puntos).

      \startlinealignment[left]
        \framed[frame=off,location=middle]{%
          \startMPcode
            fill fullcircle scaled 20 ;
          \stopMPcode}
        \hskip -2.0em % Disminuimos el espacio entre la imagen y el código
        \framed[frame=off,location=middle,align=right,width=fit]{%
          \switchtobodyfont[11.0pt]
          \starttyping[option=METAFUN]
            \startMPcode
              fill fullcircle scaled 20 ;
            \stopMPcode
          \stoptyping}
      \stoplinealignment
      
      Normalmente, se prefiere especificar el tamaño en una unidad conveniente,
      que suele ser una del sistema métrico decimal, o si conviene, del sistema
      imperial inglés.
    
    \stopsubsection

    \startsubsection[title={Operadores secundarios \color[darkgreen]{\type{xscaled}},
        \color[darkgreen]{\type{yscaled}}} y \color[darkgreen]{\type{xyscaled}}]

      En \MetaFun\ y \MetaPost, estos son operadores de transformación (un tipo
      especial de función binaria).Sirven para aplicar un escalado no uniforme a un
      objeto geométrico (path, picture o pen), estirándolo o encogiéndolo en los ejes
      X e Y de forma independiente. ¿Cómo funciona? A diferencia de scaled (que
      escala todo de forma proporcional con un solo número), xyscaled toma como
      argumento un par de coordenadas (pair), definido como (factor_x, factor_y).
      El primer número multiplica la anchura (eje X). El segundo número multiplica la
      altura (eje Y).

      Por tanto, si no queremos escalar el objeto gráfico (\type{picture}) de forma
      uniforme, tenemos las tres siguientes opciones. La primera que vamos a ver
      escala la dimensión horizontal, pero deja sin modificar la vertical:
      
      \startlinealignment[left]
        \framed[frame=off,location=middle]{%
          \startMPcode
            fill fullcircle xscaled (1cm) ;
          \stopMPcode}
        \hskip -2.0em % Disminuimos el espacio entre la imagen y el código
        \framed[frame=off,location=middle,align=right,width=fit]{%
          \switchtobodyfont[11.0pt]
          \starttyping[option=METAFUN]
            \startMPcode
              fill fullcircle xscaled 1cm ;
            \stopMPcode
          \stoptyping}
      \stoplinealignment

      La segunda escala la dimensión vertical, dejando sin modificar la horizontal:
      
      \startlinealignment[left]
        \framed[frame=off,location=middle]{%
          \startMPcode
            fill fullcircle yscaled 1cm ;
          \stopMPcode}
        \hskip -2.0em % Disminuimos el espacio entre la imagen y el código
        \framed[frame=off,location=middle,align=right,width=fit]{%
          \switchtobodyfont[11.0pt]
          \starttyping[option=METAFUN]
            \startMPcode
              fill fullcircle yscaled 1cm ;
            \stopMPcode
          \stoptyping}
      \stoplinealignment

      Por último, la tercera permite modificar las dos dimensiones a la vez.
      Los dos valores se deben situar entre paréntesis y separados por una coma:
      
      \startlinealignment[left]
        \framed[frame=off,location=middle]{%
          \startMPcode
            fill fullcircle xyscaled (2cm, 10mm) ;
          \stopMPcode}
        \hskip -3.0em % Disminuimos el espacio entre la imagen y el código
        \framed[frame=off,location=middle,align=right,width=fit]{%
          \switchtobodyfont[11.0pt]
          \starttyping[option=METAFUN]
            \startMPcode
              fill fullcircle xyscaled (2cm, 10mm) ;
            \stopMPcode
          \stoptyping}
      \stoplinealignment
      
    \stopsubsection

    \startsubsection[title=Atributo \color[darkgreen]{\type{withcolor}}]
      Es un modificador. No genera objetos, sino que modifica las propiedades de
      un comando de dibujo (como \type{draw} o \type{fill}). Representaremos
      una elipse de color rojo:
      
      \startlinealignment[left]
        \framed[frame=off,location=middle]{%
          \startMPcode
            fill fullcircle xyscaled (2cm, 1cm) withcolor "red" ;
          \stopMPcode}
        \hskip -4.0em % Disminuimos el espacio entre la imagen y el código
        \framed[frame=off,location=middle,align=right,width=fit]{%
          \switchtobodyfont[11.0pt]
          \starttyping[option=METAFUN]
            \startMPcode
              fill fullcircle xyscaled (2cm, 1cm) withcolor "red" ;
            \stopMPcode
          \stoptyping}
      \stoplinealignment

    \stopsubsection

  \stopsection
  

  \stopsection
 

\stopchapter

\stopcomponent




%%% Local Variables:
%%% coding: utf-8
%%% mode: context
%%% ConTeXt-Mark-version: "IV"
%%% TeX-engine: default
%%% TeX-master: "../ctx-tutorialLuaMetaFun.tex"
%%% End:
